\chapter{第5套试卷}

\begin{center}
\textbf{FPGA与VHDL 基础试卷（五）}

（满分：100分 \hspace{2em} 考试时间：120分钟）
\end{center}

% ============================================
\section{一、选择题（共6题，每题3分，共18分）}
% ============================================

\begin{enumerate}[label=\arabic*., leftmargin=*]

\item 下列关于CPLD与FPGA结构的说法，正确的是（~~~）。

\vspace{0.3cm}

\begin{tabular}{@{}lp{0.40\textwidth}@{\qquad}lp{0.40\textwidth}@{}}
A. & CPLD基于查找表结构，FPGA基于乘积项结构 &
B. & CPLD基于乘积项结构，适合组合逻辑；FPGA基于查找表结构，适合时序逻辑 \\
C. & CPLD和FPGA均基于查找表结构 &
D. & CPLD的集成度通常远高于FPGA \\
\end{tabular}

\vspace{0.8cm}

\item FPGA中的查找表（LUT）实质上是一个（~~~）。

\vspace{0.3cm}

\begin{tabular}{@{}lp{0.40\textwidth}@{\qquad}lp{0.40\textwidth}@{}}
A. & 由D触发器构成的阵列 &
B. & 由SRAM构成的小容量存储器 \\
C. & 专用的硬件乘法器 &
D. & 锁相环（PLL）模块 \\
\end{tabular}

\vspace{0.8cm}

\item 在VHDL中，以下关于SIGNAL（信号）与VARIABLE（变量）的描述，\textbf{错误}的是（~~~）。

\vspace{0.3cm}

\begin{tabular}{@{}lp{0.40\textwidth}@{\qquad}lp{0.40\textwidth}@{}}
A. & 信号用于进程间通信，变量只能在进程内部使用 &
B. & 变量的赋值立即生效，信号赋值在进程挂起时才更新 \\
C. & 信号使用\verb|<=|赋值，变量使用\verb|:=|赋值 &
D. & 信号只能在进程内声明，变量可以在结构体中声明 \\
\end{tabular}

\vspace{0.8cm}

\item IEEE.STD\_LOGIC\_1164标准中，STD\_LOGIC数据类型一共包含多少种逻辑值？（~~~）

\vspace{0.3cm}

\begin{tabular}{@{}lp{0.40\textwidth}@{\qquad}lp{0.40\textwidth}@{}}
A. & 2种（'0'和'1'） &
B. & 4种（'0'，'1'，'Z'，'X'） \\
C. & 7种 &
D. & 9种 \\
\end{tabular}

\vspace{0.8cm}

\item 下列关于VHDL中PROCESS（进程）语句的说法，正确的是（~~~）。

\vspace{0.3cm}

\begin{tabular}{@{}lp{0.40\textwidth}@{\qquad}lp{0.40\textwidth}@{}}
A. & 一个结构体（ARCHITECTURE）中最多只能包含一个进程 &
B. & 进程的敏感信号列表可以省略，不影响综合 \\
C. & 当敏感信号列表中任一信号发生变化时，进程被激活执行 &
D. & 进程内部不能使用IF语句，只能使用CASE语句 \\
\end{tabular}

\vspace{0.8cm}

\item PAL（可编程阵列逻辑）器件的基本结构特点是（~~~）。

\vspace{0.3cm}

\begin{tabular}{@{}lp{0.40\textwidth}@{\qquad}lp{0.40\textwidth}@{}}
A. & 与阵列可编程，或阵列固定 &
B. & 与阵列固定，或阵列可编程 \\
C. & 与阵列和或阵列均可编程 &
D. & 与阵列和或阵列均固定 \\
\end{tabular}

\end{enumerate}

% ============================================
\section{二、填空题（共4题，每题3分，共12分）}
% ============================================

\begin{enumerate}[label=\arabic*., leftmargin=*]

\item FPGA的英文全称是\underline{\hspace{3cm}}，中文译为\underline{\hspace{3cm}}。

\vspace{0.5cm}

\item 在VHDL设计中，ENTITY（实体）描述电路的\underline{\hspace{2.5cm}}，ARCHITECTURE（结构体）描述电路的\underline{\hspace{2.5cm}}。

\vspace{0.5cm}

\item VHDL中，并置（拼接）运算符符号是\underline{\hspace{1.5cm}}；信号赋值符号是\underline{\hspace{1.5cm}}；变量赋值符号是\underline{\hspace{1.5cm}}。

\vspace{0.5cm}

\item 在VHDL的PROCESS语句中，当时钟上升沿到来时执行操作的典型写法是：

\verb|IF |\underline{\hspace{2cm}}\verb|'EVENT AND |\underline{\hspace{2cm}}\verb| = '1' THEN|。

\vspace{0.5cm}

\end{enumerate}

% ============================================
\section{三、程序改错题（共1题，每题5分，共5分）}
% ============================================

\textbf{题目：}以下VHDL程序意图设计一个带异步复位和使能端的8位二进制加法计数器（模256），但程序中共有\textbf{6处错误}。请找出这些错误，并写出正确的代码。

\vspace{0.3cm}

\begin{lstlisting}[caption={带错误的8位计数器程序（共6处错误）},escapeinside={(*@}{@*)}]
LIBRARY IEEE
USE IEEE.STD_LOGIC_1164.ALL;
USE IEEE.STD_LOGIC_UNSIGNED.ALL;

ENTITY counter8 IS
  PORT(
    clk      : IN  STD_LOGIC;
    rst      : IN  STD_LOGIC;
    en       : IN  STD_LOGIC;
    count    : OUT INTEGER RANGE 0 TO 255
  );
END counter8;

ARCHITECTURE behavioral OF counter8 IS
  SIGNAL cnt : STD_LOGIC_VECTOR(7 DOWNTO 0);
BEGIN
  PROCESS(clk)
  BEGIN
    IF rst = '1' THEN
      cnt <= 0;
    ELSIF clk'EVENT AND clk = '1' THEN
      IF en = '1' THEN
        cnt = cnt + 1;
      END IF;
    END PROCESS;
  END IF;

  count <= cnt(*@；@*)
END behavioral;
\end{lstlisting}

\vspace{0.3cm}

\textbf{要求：}

(1) 逐一指出每处错误的位置和错误原因；

(2) 给出修改后的正确VHDL代码。

\vspace{0.5cm}

% ============================================
\section{四、程序阅读分析题（共1题，每题5分，共5分）}
% ============================================

\textbf{题目：}阅读以下VHDL程序，回答下列问题。

\vspace{0.3cm}

\begin{lstlisting}[caption={待分析VHDL程序}]
LIBRARY IEEE;
USE IEEE.STD_LOGIC_1164.ALL;

ENTITY circuit_x IS
  PORT(
    din  : IN  STD_LOGIC_VECTOR(2 DOWNTO 0);
    en   : IN  STD_LOGIC;
    dout : OUT STD_LOGIC_VECTOR(7 DOWNTO 0)
  );
END circuit_x;

ARCHITECTURE behav OF circuit_x IS
BEGIN
  PROCESS(din, en)
  BEGIN
    IF en = '0' THEN
      dout <= (OTHERS => '0');
    ELSE
      CASE din IS
        WHEN "000"  => dout <= "00000001";
        WHEN "001"  => dout <= "00000010";
        WHEN "010"  => dout <= "00000100";
        WHEN "011"  => dout <= "00001000";
        WHEN "100"  => dout <= "00010000";
        WHEN "101"  => dout <= "00100000";
        WHEN "110"  => dout <= "01000000";
        WHEN "111"  => dout <= "10000000";
        WHEN OTHERS => dout <= "00000000";
      END CASE;
    END IF;
  END PROCESS;
END behav;
\end{lstlisting}

\vspace{0.3cm}

\textbf{问题：}

(1) 该程序实现的是什么电路？其功能是什么？（2分）

(2) 信号\verb|en|在电路中起什么作用？当\verb|en = '0'|时输出是什么？（1分）

(3) 若输入\verb|din = "011"|且\verb|en = '1'|，输出\verb|dout|的值是什么？（1分）

(4) 该电路属于组合逻辑电路还是时序逻辑电路？请说明判断依据。（1分）

\vspace{0.5cm}

% ============================================
\section{五、综合设计题（共3题，共60分）}
% ============================================

% ---------- 设计题1 ----------
\subsection{设计题1：8线--3线优先编码器设计（20分）}

设计一个8线--3线优先编码器（8-to-3 Priority Encoder），输入为$I_7 \sim I_0$（$I_7$优先级最高），输出为$Y_2 Y_1 Y_0$（二进制编码），同时包含一个有效输出标志GS（Group Signal，当至少有一个输入有效时GS=1）。

\vspace{0.3cm}

\textbf{要求：}

(1) 列出8--3优先编码器的真值表（5分）；

(2) 写出$Y_2$、$Y_1$、$Y_0$和GS的逻辑表达式（5分）；

(3) 编写该优先编码器的VHDL代码（要求使用IF-ELSIF语句描述优先级）（6分）；

(4) 画出该编码器的外部接口（符号）框图，标注所有输入输出信号（4分）。

\vspace{0.5cm}

% ---------- 设计题2 ----------
\subsection{设计题2：模60计数器设计（20分）}

设计一个同步模60计数器（计数范围0--59），可用于数字钟的秒/分计数。计数器包含时钟输入clk、异步复位rst（高电平有效）、使能输入en（高电平有效），以及两个BCD码输出：十位输出\verb|tens|（$3\sim0$）和个位输出\verb|units|（$3\sim0$）。当计数值达到59后，下一个时钟脉冲使计数值回到0。

\vspace{0.3cm}

\textbf{要求：}

(1) 画出该模60计数器的状态转换图（至少画出0--59中若干有代表性的状态迁移，并标注复位后的初始状态）（4分）；

(2) 编写完整的VHDL代码实现该计数器（8分）；

(3) 编写一个简单的测试平台（Testbench）代码，验证计数器的功能，包括复位、使能计数和模60循环（6分）；

(4) 简述该计数器在数字钟系统中的应用原理（2分）。

\vspace{0.5cm}

% ---------- 设计题3 ----------
\subsection{设计题3：``101''序列检测器（Moore状态机）设计（20分）}

设计一个Moore型状态机，用于检测串行输入序列中的``101''模式。每当时钟上升沿到来时，采样一位串行输入din，状态机检测到完整的``101''序列后，输出dout置为'1'（持续一个时钟周期），否则dout为'0'。允许检测重叠的序列（即``10101''中应检测出两次``101''）。

\vspace{0.3cm}

例如输入序列：\verb|...0 1 0 1 1 0 1...|，则检测输出为：\verb|...0 1 0 1 0 0 1...|（加粗位置输出1）。

\vspace{0.3cm}

\textbf{要求：}

(1) 画出该Moore状态机的\textbf{状态转换图}，标注每个状态的名称、含义及其输出值（6分）；

(2) 列出\textbf{状态转换表}（当前状态、输入、次态、输出）（5分）；

(3) 编写完整的VHDL代码实现该序列检测器（用双进程或三进程方式描述）（7分）；

(4) 若不希望检测重叠序列（即``10101''中只检测出一次``101''），状态转换图应如何修改？请简要说明修改方法（2分）。

\vspace{0.3cm}

{\centering
\rule{0.8\textwidth}{0.5pt}
\vspace{0.3cm}

\textbf{—— 试卷结束 ——}

\vspace{0.5cm}}

\endinput
